% =============================================================================
% paper_starter.tex -- a minimal LaTeX starter for the Polygence paper.
%
% This is a SIMPLER alternative to paper.tex (the full Pandoc-converted file).
% Use this one to learn the LaTeX syntax; paste the actual content from paper.md
% into the section bodies as you go. Compile in Overleaf -- no local install.
%
% Things to know:
%   * Lines starting with % are comments (LaTeX ignores them).
%   * \section{Foo} starts a new top-level section. \subsection{Bar} nests one
%     level deeper.
%   * Bold = \textbf{...}, italic = \textit{...}, monospace = \texttt{...}.
%   * Math goes in $...$ for inline or \[...\] for displayed.
%   * Use \cite{key} for references; the bibliography is at the end.
% =============================================================================

\documentclass[11pt]{article}

% --- Packages we need ------------------------------------------------------
\usepackage[margin=1in]{geometry}     % standard 1-inch margins
\usepackage{amsmath, amssymb}         % math symbols and environments
\usepackage{graphicx}                 % include figures with \includegraphics
\usepackage{booktabs}                 % nice table rules (\toprule, \midrule)
\usepackage{hyperref}                 % clickable links and cross-refs
\usepackage{natbib}                   % author-year citation style
\usepackage{xcolor}                   % color text if needed
\usepackage{caption}                  % nicer figure/table captions
\usepackage{microtype}                % small typography polish

% --- Title block -----------------------------------------------------------
\title{Trust Dynamics in Multi-Agent Strategic Interaction:\\
  A Simulation Study of Bayesian Reputation Systems\\
  in 8-Player Limit Texas Hold'em}
\author{Rachit Agrawal\\
  Polygence Research Project, 2025--2026}
\date{\today}

% =============================================================================
\begin{document}
\maketitle

% --- Abstract --------------------------------------------------------------
\begin{abstract}
Reputation systems that infer trustworthiness from observed behavior are
ubiquitous in online marketplaces, social networks, and financial
platforms. Yet it remains unclear whether such systems inherently reward
cooperative behavior or create exploitable dynamics that favor aggression.
We investigate this question using a controlled multi-agent simulation of
8-player Limit Texas Hold'em, where eight strategically distinct agents
interact over hundreds of thousands of hands while maintaining Bayesian
posterior beliefs about each other's types. We report a strong
anticorrelation between trust and economic performance (Pearson
$r = -0.84$): agents perceived as most trustworthy by their opponents
accumulate the least wealth, while the least-trusted aggressive agent
dominates through fold equity rather than superior play. We further
demonstrate a fundamental classification ceiling and validate
robustness across two independent implementation phases.
\end{abstract}

% =============================================================================
\section{Introduction}
\label{sec:intro}

\subsection{Motivation}
% TODO: copy the motivation paragraphs from paper.md section 1.1 here.

\subsection{Research Question}
% TODO: copy from paper.md 1.2.

\subsection{Contributions}
% TODO: copy from paper.md 1.3. Use a numbered list:
\begin{enumerate}
  \item \textbf{Trust--profit anticorrelation.} ...
  \item \textbf{Aggression dominance via avoidance.} ...
  \item \textbf{Classification ceiling.} ...
\end{enumerate}

% =============================================================================
\section{Related Work}
\label{sec:related}

% TODO: copy section 2 from paper.md. Cite using \citep{key} for parenthetical
% citations and \citet{key} for textual ones. Define the keys at the bottom of
% this file.

% =============================================================================
\section{Methodology}
\label{sec:methods}

\subsection{Game environment}
% TODO

\subsection{Agent archetypes}
% TODO -- the eight agents. A table is the natural format here:
\begin{table}[htbp]
  \centering
  \begin{tabular}{lll}
    \toprule
    Archetype & Style & Bluff rate (BR) \\
    \midrule
    Oracle    & Nash equilibrium     & 0.33 \\
    Sentinel  & Tight-aggressive     & 0.08 \\
    Firestorm & Loose-aggressive     & 0.62 \\
    Wall      & Tight-passive        & 0.04 \\
    Phantom   & Deceptive            & 0.52 \\
    Predator  & Adaptive exploiter   & 0.21 \\
    Mirror    & Imitator             & 0.09 \\
    Judge     & Conditional retaliator & 0.08 / 0.62 \\
    \bottomrule
  \end{tabular}
  \caption{The eight agent archetypes. Bluff rate is the probability of
    betting with a weak hand. Judge has two values (cooperative /
    retaliatory).}
  \label{tab:archetypes}
\end{table}

\subsection{Bayesian trust model}
% TODO. Math example:
% Each agent maintains a posterior over opponent types updated by
\[
  P(t \mid a_{1:n}) \propto P(a_n \mid t)^\lambda \cdot P(t \mid a_{1:n-1}),
\]
% where $\lambda$ is the recency-weighting parameter.

% =============================================================================
\section{Experimental Setup}
\label{sec:setup}

% TODO: 5 seeds, 10000 hands, etc. Reference Table~\ref{tab:archetypes}.

% =============================================================================
\section{Results}
\label{sec:results}

\subsection{Phase 1: Behavioral profiles}
% TODO

\subsection{Phase 1: Trust--profit anticorrelation}
% This is the headline result. Numbers from reports/metrics_scorecard.txt:
% trust-profit r = -0.837 (Phase 1, 5 seeds x 10000 hands).

\subsection{Phase 1: Firestorm dominance via fold equity}
% TODO -- 87.1\% fold equity number from reports.

\subsection{Phase 1: Classification ceiling}
% TODO -- the Sentinel/Mirror/Judge identifiability cluster.

\subsection{Phase 2: Bounded online optimization (adaptive)}
% NEW: copy from phase2/adaptive/phase2_report.md.
% Headline: r unchanged at -0.769 vs Phase 1's -0.773.
% This section replaces the OLD Phase 2 (imitation) content in paper.md.

% =============================================================================
\section{Discussion}
\label{sec:discussion}

\subsection{Why trust is costly}
% TODO

\subsection{Implications for reputation system design}
% TODO

\subsection{Limitations}
% TODO

% =============================================================================
\section{Conclusion}
\label{sec:conclusion}

% TODO: 2-3 paragraph summary.

% =============================================================================
\section*{References}
\label{sec:references}

% Simple manual references (no .bib file needed for now). Switch to BibTeX
% later when the list grows past ~10 entries.
\begin{thebibliography}{99}

\bibitem{neumann1944}
J.~von Neumann and O.~Morgenstern,
\textit{Theory of Games and Economic Behavior}.
Princeton University Press, 1944.

\bibitem{nash1950}
J.~Nash, ``Equilibrium points in n-person games,''
\textit{Proceedings of the National Academy of Sciences},
vol.~36, no.~1, pp.~48--49, 1950.

% TODO: add the rest from paper.md's reference section.

\end{thebibliography}

\end{document}
